\chapter{Statistical Learning}\label{n:stat}
\studytime{2.5 hours}

\section{Statistical Learning: the Basic Model}

\begin{dfn}{Statistical learning}
\textbf{Statistical learning} is a set of tools for estimating the relationship $f$ between input variables $X=(X_1,\dots,X_p)$ (predictors, features) and an output $Y$ (response):
\[ Y=f(X)+\varepsilon, \]
where $f$ is fixed but unknown and $\varepsilon$ is a random error with mean zero, independent of $X$.
\end{dfn}

\begin{itemize}
  \item \textbf{Why estimate $f$?} (i)~\textbf{Prediction}: $\hat Y=\hat f(X)$ (treat $\hat f$ as a black box). (ii)~\textbf{Inference}: understand \emph{how} $Y$ depends on each $X_j$.
  \item \textbf{Reducible error:} $[f(X)-\hat f(X)]^2$, can be reduced with a better model.
  \item \textbf{Irreducible error:} $\Var(\varepsilon)=\sigma^2$, cannot be removed by any model.
\end{itemize}
\[ \E[(Y-\hat Y)^2]=\underbrace{[f(X)-\hat f(X)]^2}_{\text{reducible}}+\underbrace{\sigma^2}_{\text{irreducible}} \]

\textbf{How to estimate $f$:}
\begin{itemize}
  \item \textbf{Parametric methods:} assume a form (e.g.\ linear), then estimate a few parameters. Simple, but wrong form $\Rightarrow$ poor fit.
  \item \textbf{Non-parametric methods:} no fixed form (e.g.\ k-NN, splines). Flexible, but need much data and may overfit.
\end{itemize}

\section{Supervised vs.\ Unsupervised Learning}
\begin{itemize}
  \item \textbf{Supervised:} for every $x_i$ there is a response $y_i$; fit a model to predict $y$. Regression (quantitative $y$) or classification (qualitative $y$).
  \item \textbf{Unsupervised:} only $x_i$, no $y_i$. Goal: understand relationships between variables or observations, e.g.\ \textbf{cluster analysis}.
  \item \textbf{Semi-supervised:} $y$ known for only some observations.
\end{itemize}

\section{Loss Functions, Training and Test Loss}

\begin{dfn}{Loss function}
A \textbf{loss function} $L(y,\hat y)$ measures the penalty for predicting $\hat y$ when the true value is $y$.
\end{dfn}
\begin{center}
\small
\begin{tabular}{@{}lll@{}}
\toprule
\textbf{Loss} & \textbf{Formula} & \textbf{Used for}\\
\midrule
Squared loss & $(y-\hat y)^2$ & regression (MSE)\\
Absolute loss & $|y-\hat y|$ & regression (MAE), robust to outliers\\
0--1 loss & $\mathbb 1[y\ne\hat y]$ & classification (error rate)\\
Log-loss & $-[y\ln p+(1-y)\ln(1-p)]$ & probabilistic classification\\
\bottomrule
\end{tabular}
\end{center}

\begin{dfn}{Training loss and test loss}
\textbf{Training loss} is the average loss on the data used to fit the model: $\frac1n\sum_{i=1}^nL(y_i,\hat f(x_i))$. \textbf{Test loss} is the average loss on new, unseen data. The test loss measures how well the model \textbf{generalises}.
\end{dfn}

\begin{itemize}
  \item Training loss always falls as model \textbf{flexibility} increases.
  \item Test loss is \textbf{U-shaped}: it falls, reaches a minimum, then rises (overfitting).
  \item A small training loss does \textbf{not} guarantee a small test loss.
\end{itemize}

\begin{drawit}{Training and test loss vs.\ flexibility}
\centering
\begin{tikzpicture}
\begin{axis}[width=8cm,height=4.8cm,xlabel={model flexibility},ylabel={loss},ticks=none,xmin=0,xmax=10,ymin=0,ymax=6,legend style={font=\small,at={(1.02,0.5)},anchor=west}]
\addplot[thick,mlblue,domain=0.5:10,samples=50]{4.5*exp(-0.35*x)+0.2}; \addlegendentry{training loss}
\addplot[thick,mlred,domain=0.5:10,samples=50]{4.5*exp(-0.45*x)+0.06*(x-1)^2+1}; \addlegendentry{test loss}
\addplot[dashed,gray,domain=0:10]{1}; \addlegendentry{irreducible $\sigma^2$}
\node[font=\scriptsize] at (axis cs:1.6,5.3) {underfit};
\node[font=\scriptsize] at (axis cs:8.7,5.3) {overfit};
\end{axis}
\end{tikzpicture}
\end{drawit}

\begin{numer}[: computing losses]
$y=(3,5,2)$, $\hat y=(2.5,5,4)$. Errors $=(0.5,0,-2)$.\\
MSE $=\frac{0.25+0+4}{3}=1.417$;\quad RMSE $=\sqrt{1.417}=1.190$;\quad MAE $=\frac{0.5+0+2}{3}=0.833$.
\end{numer}

\section{Risk and Its Estimation}
\begin{dfn}{Risk and empirical risk}
The \textbf{risk} of a model $f$ is its \textbf{expected loss} over the whole population: $\mathcal R(f)=\E[L(Y,f(X))]$. Because the population is unknown, we use the \textbf{empirical risk} (average loss on a sample): $\hat{\mathcal R}(f)=\frac1n\sum_iL(y_i,f(x_i))$.
\end{dfn}

\textbf{Estimating risk (test error):}
\begin{enumerate}
  \item \textbf{Validation-set approach:} split into training and validation parts. Simple, but the estimate varies a lot and uses less data.
  \item \textbf{Leave-one-out CV (LOOCV):} train on $n-1$ points, test on 1, repeat $n$ times. Low bias, but costly.
  \item \textbf{k-fold cross-validation:} split data into $k$ folds; each fold is the validation set once:
  \[ \mathrm{CV}_{(k)}=\frac1k\sum_{i=1}^kE_i \qquad (k=5\text{ or }10). \]
\end{enumerate}

\begin{drawit}{5-fold cross-validation}
\centering
\begin{tikzpicture}[font=\scriptsize]
\foreach \r in {1,...,5}{
  \node[anchor=east] at (-0.1,-\r*0.5) {Round \r};
  \foreach \c in {1,...,5}{
    \ifnum\c=\r \fill[mlorange!60] (\c*1.3-1.3,-\r*0.5-0.2) rectangle ++(1.25,0.4);
    \else \fill[mlblue!25] (\c*1.3-1.3,-\r*0.5-0.2) rectangle ++(1.25,0.4); \fi}
  \node[anchor=west] at (6.6,-\r*0.5) {$E_\r$};}
\fill[mlblue!25] (8,-1) rectangle ++(0.4,0.3); \node[anchor=west] at (8.45,-0.85) {train};
\fill[mlorange!60] (8,-1.6) rectangle ++(0.4,0.3); \node[anchor=west] at (8.45,-1.45) {validate};
\end{tikzpicture}
\end{drawit}

\begin{numer}[: 5-fold CV]
Fold errors $0.20,0.25,0.15,0.22,0.18$. $\ \mathrm{CV}_{(5)}=\frac{1.00}{5}=\mathbf{0.20}$.
\end{numer}

\section{Trade-offs in Statistical Learning}

\subsection{Bias--variance trade-off}
\begin{dfn}{Bias and variance of a model}
\textbf{Bias} is the error from approximating a complex real problem by a simpler model (systematic error). \textbf{Variance} is the amount by which $\hat f$ would change if it were trained on a different training set (sensitivity to data).
\end{dfn}
\[
\E\big[(y_0-\hat f(x_0))^2\big]=\Bias\big[\hat f(x_0)\big]^2+\Var\big[\hat f(x_0)\big]+\sigma^2
\]
\begin{itemize}
  \item More flexible model $\Rightarrow$ bias $\downarrow$, variance $\uparrow$.
  \item Simple model: high bias (underfit). Complex model: high variance (overfit).
  \item The best model balances both, at the minimum of the test-error curve.
\end{itemize}

\begin{drawit}{Bias--variance trade-off; darts board}
\centering
\begin{tikzpicture}
\begin{axis}[width=7.2cm,height=4.6cm,xlabel={complexity},ylabel={error},ticks=none,xmin=0,xmax=10,ymin=0,ymax=6,legend style={font=\scriptsize,at={(0.5,1.03)},anchor=south,legend columns=3}]
\addplot[thick,mlblue,domain=0.5:10]{5*exp(-0.45*x)}; \addlegendentry{bias$^2$}
\addplot[thick,mlorange,domain=0.5:10]{0.05*x^2}; \addlegendentry{variance}
\addplot[very thick,mlred,domain=0.5:10]{5*exp(-0.45*x)+0.05*x^2+0.5}; \addlegendentry{total}
\end{axis}
\begin{scope}[xshift=7.6cm,yshift=0.4cm]
\foreach \i/\bx/\by/\sp/\t in {0/0/1.9/0.12/{low bias, low var},1/2.2/1.9/0.5/{low bias, high var},2/0/0/0.12/{high bias, low var},3/2.2/0/0.5/{high bias, high var}}{
\begin{scope}[shift={(\i*0+\bx,\by)}]
\draw (0,0) circle(0.75); \draw (0,0) circle(0.45); \fill[mlred!60] (0,0) circle(0.15);
\node[font=\tiny,below] at (0,-0.75) {\t};
\end{scope}}
\foreach \x/\y in {0.05/1.95,-0.08/1.88,0.1/1.82,-0.02/2.0}{\fill[mlblue] (\x,\y) circle(1.3pt);}
\foreach \x/\y in {2.6/2.2,1.8/1.5,2.5/1.5,1.9/2.3}{\fill[mlblue] (\x,\y) circle(1.3pt);}
\foreach \x/\y in {0.38/0.32,0.45/0.25,0.32/0.38,0.4/0.42}{\fill[mlblue] (\x,\y) circle(1.3pt);}
\foreach \x/\y in {2.8/0.5,2.3/-0.4,2.9/-0.1,2.5/0.65}{\fill[mlblue] (\x,\y) circle(1.3pt);}
\end{scope}
\end{tikzpicture}
\end{drawit}

\subsection{Other trade-offs}
\begin{itemize}
  \item \textbf{Prediction accuracy vs.\ interpretability:} linear regression and trees are easy to explain; boosting and neural networks are accurate but black boxes.
  \item \textbf{Flexibility vs.\ data size:} flexible models need more data.
  \item \textbf{Parsimony (Occam's razor):} prefer the simplest model that explains the data.
\end{itemize}

\section{Classification Metrics}
\begin{center}
\small
\begin{tabular}{@{}c|cc@{}}
 & \textbf{Predicted +} & \textbf{Predicted $-$}\\ \hline
\textbf{Actual +} & TP & FN\\
\textbf{Actual $-$} & FP & TN\\
\end{tabular}
\end{center}
\[
\text{Accuracy}=\frac{TP+TN}{\text{all}},\quad \text{Precision}=\frac{TP}{TP+FP},\quad \text{Recall}=\frac{TP}{TP+FN},\quad F_1=\frac{2PR}{P+R}.
\]
\begin{itemize}
  \item \textbf{Precision:} of predicted positives, how many are correct. Important when FP is costly (spam folder).
  \item \textbf{Recall (sensitivity):} of actual positives, how many are found. Important when FN is costly (cancer test).
  \item Accuracy misleads on \textbf{imbalanced} data (99\% ``not fraud'' is 99\% accurate yet useless).
\end{itemize}

\begin{numer}[: confusion matrix]
100 e-mails: $TP=40$, $FN=10$, $FP=5$, $TN=45$. Accuracy $=0.85$; Precision $=40/45=0.889$; Recall $=40/50=0.80$; $F_1=\frac{2(0.889)(0.8)}{1.689}=0.842$.
\end{numer}

\section{Sampling Distribution of an Estimator}

\begin{dfn}{Estimator and sampling distribution}
An \textbf{estimator} $\hat\theta$ is a rule (statistic) that computes an estimate of a population parameter $\theta$ from a sample (e.g.\ the sample mean $\bar X$ for $\mu$). The \textbf{sampling distribution} of $\hat\theta$ is the probability distribution of its values over all possible samples of the same size.
\end{dfn}

\begin{itemize}
  \item \textbf{Bias:} $\Bias(\hat\theta)=\E[\hat\theta]-\theta$; unbiased if $0$. $\ \mathrm{MSE}(\hat\theta)=\Var(\hat\theta)+\Bias^2$.
  \item \textbf{Standard error:} standard deviation of the sampling distribution. For the mean: $\mathrm{SE}(\bar X)=\sigma/\sqrt n$.
  \item \textbf{Central Limit Theorem:} for large $n$ ($\ge30$), $\bar X\approx\mathcal N(\mu,\sigma^2/n)$, whatever the population shape.
  \item \textbf{95\% confidence interval:} $\bar x\pm1.96\,s/\sqrt n$.
  \item \textbf{Sample variance} $s^2=\frac{1}{n-1}\sum(x_i-\bar x)^2$ is unbiased; dividing by $n$ is biased (too small).
  \item \textbf{Bootstrap:} estimate SE by resampling the data with replacement many times.
  \item 4 times more data $\Rightarrow$ SE halves.
\end{itemize}

\begin{numer}[: standard error and CI]
$\bar x=70$, $\sigma=10$, $n=25$. $\ \mathrm{SE}=10/\sqrt{25}=2$. 95\% CI $=70\pm1.96(2)=70\pm3.92=(66.08,\,73.92)$.
\end{numer}

\begin{imp}
Most asked: bias--variance trade-off (formula + diagram), training vs.\ test loss (U-curve), k-fold cross-validation, standard error and CLT, precision/recall numerical.
\end{imp}

\begin{iq}
\item \twom What is statistical learning? Write $Y=f(X)+\varepsilon$ and explain.
\item \twom Differentiate reducible and irreducible error.
\item \twom Define risk and empirical risk.
\item \twom What is the standard error of the mean?
\item \fivem Compare parametric and non-parametric methods.
\item \fivem Explain k-fold cross-validation with a diagram.
\item \fivem Explain training loss and test loss with a graph.
\item \fivem State the central limit theorem and its importance.
\item \tenm Explain the bias--variance trade-off with formula and diagram.
\item \tenm Explain the sampling distribution of an estimator, bias, standard error and confidence interval with an example.
\end{iq}
